% ==============================================================================
% 09_criteria6_studentsupport.tex — Criteria 6: Student Support Service
% Score: 2
% ==============================================================================

\criterion{6}{บริการสนับสนุนผู้เรียน}{Student Support Service}

% ------------------------------------------------------------------------------
\criterionsub{6.1 The student intake policy, admission criteria, and admission procedures to the programme are shown to be clearly defined, communicated, published, and up-to-date.}

หลักสูตรวิศวกรรมศาสตรมหาบัณฑิต สาขาวิชาวิศวกรรมคอมพิวเตอร์และปัญญาประดิษฐ์ (CE\&AI) กำหนดนโยบายและเกณฑ์การรับเข้าตาม\textbf{ข้อบังคับมหาวิทยาลัยราชภัฏอุตรดิตถ์ ว่าด้วยการจัดการศึกษาระดับบัณฑิตศึกษา พ.ศ.~2566} และประกาศการรับสมัครของบัณฑิตวิทยาลัยในแต่ละปีการศึกษา โดยรับผู้สำเร็จการศึกษาระดับปริญญาตรีในสาขาวิศวกรรมคอมพิวเตอร์ วิทยาการคอมพิวเตอร์ เทคโนโลยีสารสนเทศ หรือสาขาอื่นที่มีความเชี่ยวชาญใกล้เคียงตามเกณฑ์ที่มหาวิทยาลัยกำหนด

\textbf{การเผยแพร่และการสื่อสาร} หลักสูตรกำหนดให้มีการเผยแพร่ข้อมูลการรับเข้าให้ครบถ้วน ทันสมัย และเข้าถึงได้หลายช่องทาง ดังสรุปในตารางที่ \ref{tbl:c6_admission_channels}

\begin{table}[H]
  \centering
  \caption{ช่องทางเผยแพร่นโยบายและขั้นตอนการรับเข้าหลักสูตร CE\&AI}
  \label{tbl:c6_admission_channels}
  \begingroup
  \small
  \setlength{\tabcolsep}{4pt}
  \renewcommand{\arraystretch}{1.2}
  \begin{tabularx}{\linewidth}{|>{\raggedright\arraybackslash}p{3.2cm}|Y|>{\centering\arraybackslash}p{2.6cm}|}
    \hline
    \rowcolor{gray!15}
    \textbf{ช่องทาง} & \textbf{เนื้อหาที่เผยแพร่} & \textbf{ความถี่การทบทวน} \\
    \hline
    เว็บไซต์มหาวิทยาลัย/บัณฑิตวิทยาลัย & ประกาศรับสมัคร เกณฑ์คุณสมบัติ ค่าธรรมเนียม ปฏิทิน & ทุกรอบรับเข้า \\
    \hline
    เว็บไซต์คณะเทคโนโลยีอุตสาหกรรม & โครงสร้างหลักสูตร PLO/CLO แนวทางอาชีพ & อย่างน้อยปีละ 1 ครั้ง \\
    \hline
    เพจ Facebook หลักสูตร CE\&AI & ประชาสัมพันธ์ Q\&A กิจกรรมเปิดบ้าน & ตลอดปีการศึกษา \\
    \hline
    เอกสารรับสมัคร/แผ่นพับ & สรุปเกณฑ์และขั้นตอนแบบพกพา & ปรับตามรอบรับสมัคร \\
    \hline
    วันปฐมนิเทศนักศึกษาใหม่ & ชี้แจงระเบียบ สิทธินักศึกษา ช่องทางรับความช่วยเหลือ & รายภาคการศึกษาแรกที่มีนักศึกษา \\
    \hline
  \end{tabularx}
  \endgroup
\end{table}

คาดการณ์รับเข้าเป้าหมายต่อปีการศึกษา \textbf{15--20 คน} (แบ่งแผน 1 / แผน 2 ตามโครงสร้างหลักสูตร) โดยในปีการศึกษา 2568 มีการรับนักศึกษาเข้าศึกษาจริงจำนวน \textbf{4 คน} ใน\textbf{ภาคการศึกษาที่ 2} และหลักสูตรกำหนดให้การรับนักศึกษาเข้าใหม่ในปีการศึกษา 2569 อยู่ใน\textbf{ภาคการศึกษาที่ 1} ตามตารางด้านล่าง

\begin{table}[H]
  \centering
  \caption{จำนวนนักศึกษารับเข้า (คน) ของหลักสูตร}
  \begin{tabularx}{\linewidth}{|>{\raggedright\arraybackslash}p{5.5cm}|Z|Z|Z|}
    \hline
    \rowcolor{gray!15}
    \textbf{ปีการศึกษา} & \textbf{แผน 1} & \textbf{แผน 2} & \textbf{รวม} \\
    \hline
    2568 (เป้าหมาย) & 5--10 & 10--15 & 15--20 \\
    \hline
    2568 (รับเข้าจริง ภาค 2) & --- & --- & 4 \\
    \hline
    2569 (ภาค 1 — อยู่ระหว่างรับเข้า) & \multicolumn{3}{c|}{สัมภาษณ์ผู้สมัคร 8 คน (ยังไม่ประกาศผลการคัดเลือก)} \\
    \hline
  \end{tabularx}
\end{table}

\begin{table}[H]
  \centering
  \caption{จำนวนนักศึกษาคงเหลือ (คน) ของหลักสูตร}
  \begin{tabularx}{\linewidth}{|Z|Z|Z|Z|}
    \hline
    \rowcolor{gray!15}
    \textbf{ปีการศึกษา} & \textbf{แผน 1} & \textbf{แผน 2} & \textbf{รวม} \\
    \hline
    2568 (ณ ช่วงมีการเรียนการสอน ภาค 2) & --- & --- & 4 \\
    \hline
  \end{tabularx}
\end{table}

\noindent\textit{หมายเหตุ: คอลัมน์แผน 1/แผน 2 ของจำนวนรับเข้าจริงจะสรุปจากระบบ REG เมื่อจัดประเภทแผนการเรียนได้ครบ}

\marknote{ควรอัปเดตแยกจำนวนตามแผน 1 และแผน 2 จากระบบ REG และแนบประกาศรับสมัครฉบับเป็นทางการเป็น AUNQA-6-2}

% ------------------------------------------------------------------------------
\criterionsub{6.2 Both short-term and long-term planning of academic and non-academic support services are shown to be carried out to ensure sufficiency and quality of support services for teaching, research, and community service.}

มหาวิทยาลัยและคณะเทคโนโลยีอุตสาหกรรมจัดบริการสนับสนุนการเรียนรู้และบริการนอกชั้นเรียนครบวงจร หลักสูตร CE\&AI ได้วางแผนให้สอดคล้องกับวงจรการเปิดหลักสูตรใหม่ โดยแยก\textbf{ระยะสั้น} (เตรียมความพร้อมก่อนและระหว่างเปิดรับรุ่นแรก) และ\textbf{ระยะยาว} (หลังมีบัณฑิตรุ่นแรกและข้อมูลย้อนหลัง) ตามตารางที่ \ref{tbl:c6_support_plan}

\begin{table}[H]
  \centering
  \caption{แผนบริการสนับสนุนนักศึกษา (ระยะสั้นและระยะยาว)}
  \label{tbl:c6_support_plan}
  \begingroup
  \small
  \setlength{\tabcolsep}{3pt}
  \renewcommand{\arraystretch}{1.25}
  \begin{tabularx}{\linewidth}{|>{\raggedright\arraybackslash}p{2.8cm}|Y|>{\raggedright\arraybackslash}p{3.0cm}|>{\raggedright\arraybackslash}p{2.6cm}|}
    \hline
    \rowcolor{gray!15}
    \textbf{กลุ่มบริการ} & \textbf{ระยะสั้น (2568)} & \textbf{ระยะยาว (2569--2572)} & \textbf{ผู้รับผิดชอบหลัก} \\
    \hline
    Academic (ห้องสมุด/ฐานข้อมูล) & ยืนยันสิทธิการเข้าถึง EZproxy/VPN สำหรับนักศึกษาใหม่ จัดอบรมการสืบค้นวรรณกรรม & ทบทวนสิทธิฐานข้อมูลรายปี เก็บสถิติการใช้งานรายหลักสูตร & สำนักวิทยบริการ + หลักสูตร \\
    \hline
    Academic (LMS/ห้องปฏิบัติการ) & ตั้งรายวิชาบน LMS ทุกรายวิชาเปิดสอน ตรวจสอบความพร้อมห้องปฏิบัติการ AI/ML & พัฒนาเทมเพลตรายวิชา CE\&AI ร่วมกับศูนย์คอมพิวเตอร์ & อาจารย์ผู้รับผิดชอบรายวิชา + ศูนย์คอมพิวเตอร์ \\
    \hline
    Non-academic (ทุน/ที่พัก/สุขภาพ) & ประชาสัมพันธ์ทุนของมหาวิทยาลัย ช่องทางหอพัก บริการพยาบาล & สรุปข้อมูลการใช้ทุน/ความต้องการที่พักของนักศึกษา CE\&AI รายปี & ฝ่ายกิจการนักศึกษา \\
    \hline
    การที่ปรึกษา/ติดตามความก้าวหน้า & มอบหมายอาจารย์ที่ปรึกษาประจำตัวและที่ปรึกษาวิทยานิพนธ์ตามแผน & พัฒนาแดชบอร์ดติดตามความก้าวหน้า (ร่วมกับ REG) & ประธานหลักสูตร/งานบัณฑิตศึกษา \\
    \hline
    Employability & จัดสัมมนาเชิงวิชาการเชิญวิทยากรภายนอก & ขยายเครือข่ายสถานประกอบการ/ศิษย์เก่า สำหรับแนวทางอาชีพ & หลักสูตร + คณะ \\
    \hline
  \end{tabularx}
  \endgroup
\end{table}

นอกจากนี้ บริการที่มีอยู่แล้วในระดับมหาวิทยาลัย ได้แก่
ประกอบด้วย \textbf{Academic support} จากสำนักวิทยบริการ (ห้องสมุดกลาง) และศูนย์คอมพิวเตอร์ ซึ่งมีฐานข้อมูล IEEE Xplore, Springer Nature Link, ScienceDirect, Engineering Source, Emerald ระบบ LMS และซอฟต์แวร์/เครื่องมือที่สนับสนุนการเรียนการสอนและการวิจัยด้าน AI/ML; \textbf{Non-academic support} จากฝ่ายกิจการนักศึกษาในด้านทุนการศึกษา ที่ปรึกษาด้านการเงิน หอพัก และบริการสุขภาพ; และ \textbf{การลงทุนระยะยาว} ตามแผนปรับปรุงห้องปฏิบัติการ AI ของคณะและแผนพัฒนาทรัพยากรดิจิทัลของสำนักวิทยบริการให้สอดคล้องกับแผนมหาวิทยาลัย

% ------------------------------------------------------------------------------
\criterionsub{6.3 An adequate system is shown to exist for student progress, academic performance, and workload monitoring. Student progress, academic performance, and workload are shown to be systematically recorded and monitored. Feedback to students and corrective actions are made where necessary.}

ระบบติดตามความก้าวหน้าและภาระงานของนักศึกษาประกอบด้วย
ได้แก่ระบบสารสนเทศการลงทะเบียน (REG) ของกองบริการการศึกษาซึ่งบันทึกผลการเรียน หน่วยกิต และสถานะทางทะเบียน อาจารย์ที่ปรึกษาประจำตัวซึ่งติดตาม transcript และแผนการเรียนทุกภาคการศึกษา อาจารย์ที่ปรึกษาวิทยานิพนธ์/สารนิพนธ์ซึ่งพบผู้เรียนอย่างน้อย 2 ครั้งต่อเดือนตามแนวปฏิบัติของหลักสูตร และรายงานความก้าวหน้าวิทยานิพนธ์/สารนิพนธ์ในทุกปลายภาคการศึกษา ซึ่งใช้เป็นฐานในการกำหนดมาตรการแก้ไขเมื่อความก้าวหน้าไม่เป็นไปตามแผน

\textbf{จุดตรวจสอบที่วางแผนไว้} (ใช้เป็นแม่แบบเมื่อมีนักศึกษา) สรุปในตารางที่ \ref{tbl:c6_progress_checkpoints}

\begin{table}[H]
  \centering
  \caption{จุดตรวจสอบความก้าวหน้าและการให้ข้อเสนอแนะแก่นักศึกษา}
  \label{tbl:c6_progress_checkpoints}
  \begingroup
  \small
  \begin{tabularx}{\linewidth}{|>{\raggedright\arraybackslash}p{3.0cm}|Y|>{\raggedright\arraybackslash}p{3.4cm}|}
    \hline
    \rowcolor{gray!15}
    \textbf{ช่วงเวลา} & \textbf{สิ่งที่ติดตาม} & \textbf{การตอบสนอง/แก้ไข} \\
    \hline
    กลางภาคการศึกษา & ผลคะแนนเชิงตัวชี้วัดในรายวิชา ภาระงานหน่วยกิต & แจ้งเตือนนักศึกษาและอาจารย์ที่ปรึกษาเมื่อมีความเสี่ยงต่ำกว่าเกณฑ์ \\
    \hline
    ปลายภาคการศึกษา & สรุปผลการเรียน แผนการเรียนถัดไป & ปรับแผนการเรียนร่วมกับอาจารย์ที่ปรึกษา \\
    \hline
    รายงานความก้าวหน้าวิจัย & บทที่ทำได้ งานทดลอง/ผลการทดลอง & มติคณะกรรมการบัณฑิตศึกษา/ข้อเสนอแนะจากที่ปรึกษา \\
    \hline
  \end{tabularx}
  \endgroup
\end{table}

% ------------------------------------------------------------------------------
\criterionsub{6.4 Co-curricular activities, student competition, and other student support services are shown to be available to improve learning experience and employability.}

หลักสูตรดำเนินกิจกรรมเสริมหลักสูตรเพื่อเสริมประสบการณ์การเรียนรู้และความพร้อมสู่การประกอบอาชีพ โดยเผยแพร่ข่าวสารและบันทึกกิจกรรมผ่าน\textbf{เพจ Facebook หลักสูตร CE\&AI} (\texttt{facebook.com/profile.php?id=61577236346855}) ซึ่งเป็นช่องทางหลักในการสื่อสารกับนักศึกษาปัจจุบัน ผู้สมัคร และผู้มีส่วนได้ส่วนเสีย ณ ปีการศึกษา 2568 เพจมีผู้ติดตาม \textbf{491 คน} และมีการโพสต์กิจกรรมรวม \textbf{99 โพสต์} ครอบคลุม 11 ใน 12 เดือน (ร้อยละ 91.7)

\begin{table}[H]
  \centering
  \caption{กิจกรรมเสริมหลักสูตร CE\&AI ตามช่วงเวลา (ปีการศึกษา 2568)}
  \label{tbl:c6_cocurricular}
  \begingroup
  \small
  \begin{tabularx}{\linewidth}{|>{\raggedright\arraybackslash}p{3.4cm}|Y|>{\centering\arraybackslash}p{2.2cm}|}
    \hline
    \rowcolor{gray!15}
    \textbf{ช่วงเวลา / กิจกรรม} & \textbf{รายละเอียด} & \textbf{หลักฐาน} \\
    \hline
    มิ.ย. 68 — เปิดตัวหลักสูตร & เปิดตัวเพจ Facebook ประชาสัมพันธ์รับสมัครนักศึกษารุ่นแรก (16 โพสต์) & Facebook Page \\
    \hline
    ต.ค.–พ.ย. 68 — กิจกรรมวิชาการ & กิจกรรมปลายภาค 1/68 ร่วมกับศูนย์ส่งเสริมบัณฑิตศึกษา ม.ร.อ. (22 โพสต์) & Facebook Page \\
    \hline
    ม.ค.–ก.พ. 69 — เปิดภาคและรับสมัคร & กิจกรรมเข้มข้นสุด ร่วมกับคณะเทคโนโลยีอุตสาหกรรม ประชาสัมพันธ์รับสมัคร 1/69 (29 โพสต์) & Facebook Page \\
    \hline
    มี.ค. 69 — การบริการวิชาการชุมชน & ร่วมกับสถาบันการอาชีวศึกษาภาคเหนือ 3, โรงเรียนเทศบาลท่าอิฐ, กลุ่มการพยาบาล รพ.อุตรดิตถ์ (12 โพสต์) & Facebook Page \\
    \hline
    ตลอดปี — การแข่งขันและ workshop & แผนส่งนักศึกษาร่วมแข่งขันด้าน AI/นวัตกรรม และ workshop MLOps, Edge AI & AUNQA-6-5 (ค้างรวบรวม) \\
    \hline
  \end{tabularx}
  \endgroup
\end{table}

กิจกรรมที่โดดเด่น ได้แก่ การเชื่อมโยงกับหน่วยงานภายนอก \textbf{5 แห่ง} ผ่านการ cross-share เนื้อหา ประกอบด้วยคณะเทคโนโลยีอุตสาหกรรม ม.ร.อ. สถาบันการอาชีวศึกษาภาคเหนือ 3 โรงเรียนเทศบาลท่าอิฐ กลุ่มการพยาบาลโรงพยาบาลอุตรดิตถ์ และศูนย์ส่งเสริมบัณฑิตศึกษา ม.ร.อ.\ ซึ่งแสดงถึงการเสริมสร้าง learning experience ผ่านความร่วมมือกับชุมชนและภาคส่วนที่เกี่ยวข้อง

\marknote{AUNQA-6-5 = ภาพหน้าจอ/รายการโพสต์กิจกรรมจาก Facebook Page ปีการศึกษา 2568 (99 โพสต์, 11 เดือน) — ให้รวบรวมเป็นไฟล์ PDF/ตารางสรุป}

% ------------------------------------------------------------------------------
\criterionsub{6.5 The competences of the support staff rendering student services are shown to be identified for recruitment and deployment. These competences are shown to be evaluated to ensure their continued relevance to stakeholders.}

บุคลากรสายสนับสนุนที่เกี่ยวข้องกับการบริการนักศึกษาในระดับมหาวิทยาลัยและคณะ ได้แก่ กลุ่มงานกิจการนักศึกษา สำนักวิทยบริการ ศูนย์คอมพิวเตอร์ งานบัณฑิตศึกษา และงานการเงิน มีการกำหนด\textbf{สมรรถนะตามตำแหน่ง (Position Profile)} การสรรหาและการวางตัวตามภารกิจ และมีการ\textbf{ประเมินสมรรถนะรายปี}โดยกองบริหารงานบุคคล รวมถึงการประเมินความพึงพอใจของผู้รับบริการจากนักศึกษาและอาจารย์ เพื่อนำผลไปปรับปรุงการบริการ

หลักสูตร CE\&AI จะประสานกับหน่วยงานเหล่านี้เป็นประจำผ่านประธานหลักสูตรและงานบัณฑิตศึกษาของคณะ เพื่อให้ความต้องการของนักศึกษาระดับบัณฑิตศึกษาถูกสื่อสารไปยังผู้ให้บริการ

% ------------------------------------------------------------------------------
\criterionsub{6.6 Student support services are shown to be subjected to evaluation, benchmarking, and enhancement.}

บริการสนับสนุนผู้เรียนของมหาวิทยาลัยถูกประเมินผ่าน
แบบสำรวจความพึงพอใจของนักศึกษาประจำปี รายงานประเมินคุณภาพการศึกษาระดับมหาวิทยาลัย/คณะ (SAR) และการเปรียบเทียบแนวปฏิบัติ (benchmark) กับมหาวิทยาลัยในกลุ่มราชภัฏและกลุ่มพลิกโฉม โดยเฉพาะด้านบริการดิจิทัลและการสนับสนุนนักศึกษาบัณฑิตศึกษา เพื่อนำผลไปปรับปรุงบริการอย่างต่อเนื่อง

\textbf{แผนปรับปรุงเฉพาะหลักสูตร} เมื่อมีนักศึกษาแล้ว หลักสูตรจะเก็บข้อมูลความพึงพอใจเฉพาะกลุ่ม CE\&AI (เช่น การใช้ห้องปฏิบัติการ การสนับสนุนที่ปรึกษา การใช้ฐานข้อมูล) อย่างน้อยปีละ 1 ครั้ง และสรุปเป็นรายงานปรับปรุงต่อคณะกรรมการหลักสูตร

\marknote{ควรจัดเก็บผลการประเมินความพึงพอใจของนักศึกษาบัณฑิตศึกษา CE\&AI อย่างน้อยปีละครั้ง และสรุปเป็น AUNQA-6-4 เมื่อมีรอบสำรวจครบ}

\begin{selfassessmentbox}{2}{มีนโยบายรับเข้า ช่องทางเผยแพร่ และแผนบริการสนับสนุนระยะสั้น--ระยะยาวครบถ้วน มีระบบติดตามความก้าวหน้าและแผนกิจกรรมเสริมชัดเจน และเริ่มมีนักศึกษาในหลักสูตร (ปีการศึกษา 2568 จำนวน 4 คน ภาค 2) อย่างไรก็ตาม ยังต้องพัฒนาแบบสำรวจความพึงพอใจเฉพาะหลักสูตรและข้อมูล benchmark ให้ครบถ้วน}
\end{selfassessmentbox}

% ------------------------------------------------------------------------------
\subsection*{รายการหลักฐาน}

หลักฐานประกอบ Criterion~6 ประกอบด้วย AUNQA-6-1 ข้อบังคับมหาวิทยาลัยฯ ว่าด้วยการจัดการศึกษาระดับบัณฑิตศึกษา พ.ศ.~2566, AUNQA-6-2 ประกาศการรับสมัครนักศึกษาบัณฑิตศึกษา ปีการศึกษา 2568, AUNQA-6-3 รายชื่อฐานข้อมูลและบริการของสำนักวิทยบริการ ม.ราชภัฏอุตรดิตถ์, AUNQA-6-4 รายงานความพึงพอใจของนักศึกษาบัณฑิตศึกษา/แบบสำรวจเฉพาะหลักสูตร (อยู่ระหว่างพัฒนาเครื่องมือและรอข้อมูลหลังเปิดรับนักศึกษา) และ AUNQA-6-5 ปฏิทินกิจกรรมเสริมและหลักฐานการจัดกิจกรรม (เติมเมื่อเริ่มดำเนินการ)
